
% ДЛЯ КА ВВ и КА СВ! КА СВ отличие только во входах на включение выключателя в схеме В используются xlsx VV_SwitchDevice и SV_SwitchDevice
\needspace{3\baselineskip}
\color{uniblue}{\subsection{Регулирование напряжения параллельно работающих трансформаторов}}
\color{black}

\begin{enumerate}[label=\arabic{section}.\arabic{subsection}.\arabic{enumi}, labelsep=4pt, leftmargin=0pt, itemindent=57pt, itemsep=0pt, parsep=5pt] % если что можно крутить itemsep

%%%%%%%%%%%%%%%%%%%%%%%%%%%%%%%%%%%%%%%%%%%%%%%%%%%%%%%%%%%%%%%%%%%%% ОБЩИЕ СВЕДЕНИЯ %%%%%%%%%%%%%%%%%%%%%%%%%%%%%%%%%%%%%%%%%%%%%%%%%%%%%%%%%%%%%%


\item Регулирование напряжения параллельно работающих трансформаторов
   
\item Общие сведения\\
    Для двух параллельно работающих трансформаторов  в устройстве предусмотрена возможность организации регулирования напряжения по методу следования командам ведущего в автоматическом и ручном режиме управления. Для этого привода РПН трансформаторов должны управляться устройствами <<ЮНИТ-М300-АРНТ>>. В процессе регулирования контролируется отсутствие расхождения положений  РПН трансформаторов.
\item Коммуникация между устройствами
    \begin{enumerate}[label=\arabic{section}.\arabic{subsection}.\arabic{enumi}, labelsep=4pt, leftmargin=0pt, itemindent=57pt, itemsep=0pt, parsep=5pt]
        \item Обмен дискретными сигналами между устройствами в режиме регулирования параллельно работающих трансформаторов можно организовать следующими способами:
            \begin{itemize}
                \item	контрольным кабелем и дискретными входами и выходными реле терминалов (исполнение I);
                \item	GOOSE сообщениями (исполнения I и II).
            \end{itemize}
        \item Входные сигналы для исполнения I, передаваемые посредством контрольного кабеля (необходимый минимум для организации регулирования параллельно работающих трансформаторов):
            \begin{itemize}
                \item	команда <<Прибавить>> от устройства регулирования параллельно работающего трансформатора;
                \item	команда <<Убавить>> от устройства регулирования параллельно работающего трансформатора;
                \item сигнал <<Идет переключение>> от устройства регулирования параллельно работающего трансформатора;
                \item	устройство регулирования параллельно работающего трансформатора заблокировано.
            \end{itemize} 
        \item 
        В приложении Ж приведена схема функциональных взаимосвязей между устройствами при регулировании напряжения двух включенных на параллельную работу трансформаторов. 
        \item Дополнительные входные сигналы для исполнений I и II, передаваемые посредством GOOSE сообщения:
            \begin{itemize}
                \item сигнал готовности к режиму регулирования параллельно работающего трансформатора;
                \item	устройство регулирования параллельно работающего трансформатора в режиме автоматического регулирования;
                \item ненормальный режим на секции параллельно работающего трансформатора;
                \item параллельно работающий трансформатор не имеет связи с шинами;
                \item устройство регулирования параллельно работающего трансформатора в режиме <<Ведущий>>;
                \item устройство регулирования параллельно работающего трансформатора в режиме <<Принудительно ведущий>>;
                \item текущее положение РПН параллельно работающего трансформатора.
            \end{itemize} 
    \end{enumerate}

\item Настройка устройств
    \begin{enumerate}[label=\arabic{section}.\arabic{subsection}.\arabic{enumi}, labelsep=4pt, leftmargin=0pt, itemindent=57pt, itemsep=0pt, parsep=5pt]
        \item 
        Кроме коммуникации между устройствами для организации работы параллельно включенных трансформаторов необходимо в устройствах, управляющих приводами РПН каждого трансформатора, выполнить определенные настройки.
        \item 
        Необходимо с помощью программного переключателя SGF1 (см. рисунок 2.14) определить для каждого из устройств роль в  режиме параллельного регулирования – ведущий/ведомый.
        \item
        Если предполагается автоматическое регулирование напряжения, в устройстве, которое является <<Ведущим>>, следует задать автоматический режим управления.
        \item
        Управление приводами параллельно работающих трансформаторов в ручном режиме возможно с любого устройства. При этом устройство, которое выбрано для оперативного управления в ручном режиме, определяется термином <<Принудительно ведущий>>, а устройство, параллельно работающего трансформатора, следует командам управления от <<Принудительно ведущего>>.  
    \end{enumerate}

\item Контроль рассогласования приводов
    \begin{enumerate}[label=\arabic{section}.\arabic{subsection}.\arabic{enumi}, labelsep=4pt, leftmargin=0pt, itemindent=57pt, itemsep=0pt, parsep=5pt]
        \item Каждое из устройств, управляющее приводами РПН параллельно работающих трансформаторов, может получать информацию о текущем положении переключателя смежного трансформатора двумя способами (определяется программным переключателем SGF2):
            \begin{itemize}
                \item аналоговый вход типа <<токовая петля 4(0)-20 мА>>;
                \item GOOSE сообщение.
            \end{itemize}
        \item
        Выше было отмечено, что в аппаратном обеспечении устройства предусмотрены два входа типа токовая петля 4(0)-20 мА. Первый из них задействован для подключения внешнего логометра трансформатора. Назначение второго входа  - подключение токовой петли логометра параллельно работающего трансформатора.
        \item При подключении токовой петли второго входа в качестве источника информации о положении РПН параллельно работающего трансформатора для его настройки необходимо задать следующие параметры:
            \begin{itemize}
                \item значение токовой петли соответствующее начальному положению;
                \item значение токовой петли соответствующее конечному положению;
                \item значение начального положения;
                \item значение начального положения.
            \end{itemize}
        \item 
        При расхождении положения РПН параллельно работающих трансформаторов на определенную уставками величину в течение определенного времени логикой функции формируется сигнал <<Рассогласование параллельно работающих трансформаторов>>.
    \end{enumerate}



\item 
Параметры для настройки функции <<Регулирование параллельно работающих трансформаторов>> приведены в таблице 2.7.




\small
\begin{longtable}{|>{\centering\arraybackslash}m{5.3cm}|>{\centering\arraybackslash}m{3.3cm}|>{\centering\arraybackslash}m{4.2cm}|>{\centering\arraybackslash}m{1.8cm}|>{\centering\arraybackslash}m{1cm}|}
\caption{Параметры для настройки функции <<Регулирование параллельно работающих трансформаторов>>\hfill\vspace{-0.5\baselineskip}}\label{mtzvn:tbl1}\\ 
\hline
\rowcolor{gray!30}
Параметр (Параметр на ИЧМ) & Условное обозначение на схеме & Значение/ Диапазон & Единица измерения & Шаг \\ 
\hline
\endfirsthead
%\caption{\emph{Продолжение\hfill\vspace{-0.5\baselineskip}}} \\ % Отступ обозначения таблицы от таблицы и выравнивание надписи слева
\caption*{\hspace{3pt}\emph{Продолжение таблицы \ref{mtzvn:tbl1}\hfill\vspace{-0.5\baselineskip}}} \\ % сделано по ГОСТ 2.105 п.6.8.7
\hline
\rowcolor{gray!30}
Параметр (Параметр на ИЧМ) & Условное обозначение на схеме & Значение/ Диапазон & Единица измерения & Шаг \\ 
%\hline
\endhead
%\hline
\endfoot
%\hline
\endlastfoot

%==+t1*ATCC1|_TTCC5> АРН
\multicolumn{5}{|c|}{ АРН } \\ \hline 
\centering Метод регулирования трансформаторов, включенных на параллельную работу (ПарРежим) & \centering ПарРежим & \centering 1 = Не используется\\2 = Ведущий\\3 = Ведомый & \centering -- & \centering \arraybackslash -- \\
\hline
\centering Источник данных о положении РПН параллельно работающего трансформатора (ПарТипДатчика) & \centering ПарТипДатчика & \centering 0 = Контроль рассогласования не предусмотрен\\1 = Токовая петля 4(0)-20 мА\\2 = GOOSE & \centering -- & \centering \arraybackslash -- \\
\hline
\centering Сигнал токовой петли, соответствующий конечного положения РПН параллельно работающего трансформатора (ПарIмаМакс) & \centering ПарIмаМакс & \centering 5,00 ... 100,00 & \centering мА & \centering \arraybackslash 0,01 \\
\hline
\centering Сигнал токовой петли, соответствующий начального положения РПН параллельно работающего трансформатора (ПарIмаМин) & \centering ПарIмаМин & \centering 0,00 ... 4,00 & \centering мА & \centering \arraybackslash 0,01 \\
\hline
\centering Конечное положение РПН параллельно работающего трансформатора (ПарМаксN) & \centering ПарМаксN & \centering 1 ... 40 & \centering -- & \centering \arraybackslash 1 \\
\hline
\centering Начальное положение РПН параллельно работающего трансформатора (ПарМинN) & \centering ПарМинN & \centering 1 ... 40 & \centering -- & \centering \arraybackslash 1 \\
\hline
\centering Допустимое рассогласование положений приводов РПН при регулировании тр-ров включенных на параллельную работу (МаксРазбегN) & \centering МаксРазбегN & \centering 0 ... 40 & \centering -- & \centering \arraybackslash 1 \\
\hline
\centering Допустимое время рассогласования положений РПН при регулировании параллельно работающих трансформаторов (Тразбега) & \centering Тразбега & \centering 1 ... 100000 & \centering мс & \centering \arraybackslash 1 \\
\hline
%===t1


\end{longtable}
\normalsize















% % ФОРМУЛА №1
% \begin{equation*}
% \begin{bmatrix}
% I_{A,n*}\\ 
% I_{B,n*}\\ 
% I_{C,n*}
% \end{bmatrix}= k_{AM,n}\times \begin{bmatrix}
% M_{11,n} & M_{12,n} & M_{13,n}\\ 
% M_{21,n} & M_{22,n} & M_{23,n}\\ 
% M_{31,n} & M_{32,n} & M_{33,n}
% \end{bmatrix}\times \begin{bmatrix}
% I_{A,n}\\ 
% I_{B,n}\\ 
% I_{C,n}
% \end{bmatrix},
% \end{equation*}
% \begin{eqexpl}[25mm]
% \item{$n$}индекс, обозначающий сторону (плечо) ДЗТ;
% \item{$k_{AM,n}$}коэффициент амплитудного выравнивания;
% \item{$I_{A,n},I_{B,n},I_{C,n}$}измеренные фазные токи стороны $n$;
% \item{$M_{11,n}\dots M_{33,n} $}коэффициенты матрицы компенсации фазового сдвига в соответствие  с группой соединения обмоток трансформатора.
% \end{eqexpl}


%%%%%%%%%%%%%%%%%%%%%%%%%%%%%%%%%%%%%%%%%%%%%%%%%%%%%%%%%%% КОНЕЦ Выкатной элемент (ВЭ) %%%%%%%%%%%%%%%%%%%%%%%%%%%%%%%%%%%%%%%%%%%%%%%%%%%%%%%%%%


\end{enumerate}